\begin{abstract}
\textbf{EN}---Many modern scientific and machine learning workloads exhibit inherently sparse and irregular structure. However, traditional supercomputers and datacenters are optimized for dense, regular computation. As a result, sparse workloads use only a fraction of available system resources, leaving significant performance untapped. Although prior solutions have been proposed, they typically target narrow classes of sparse tensor algebra operations, do not fully leverage modern vector compute cores, or leave much of the programmability burden to the user, limiting widespread adoption. To address this gap, this thesis presents a hardware--software co-design that enables the efficient execution of general sparse tensor algebra operations at scale. We begin by analyzing the architectural implications of sparsity across a broad range of scientific and machine learning applications. Our analysis reveals that sparse tensor traversal and sparse tensor merging are the two dominant bottlenecks on conventional architectures, often stalling execution for up to 99\% of the time. Building on these insights, we introduce the \textit{Tensor Marshaling Unit} (TMU), a specialized architectural unit that accelerates these operations. We demonstrate that integrating the TMU into a multicore CPU improves performance by an order of magnitude, outperforming state-of-the-art GPUs in both performance and energy efficiency. To deliver these performance gains without exposing the complexity of DAE programming to end users, we also present Ember, a compiler that lowers sparse machine learning operations to \textit{Decoupled Access--Execute} (DAE) code for architectures such as TMU--CPU systems. Ember is implemented in MLIR, enabling natural integration with PyTorch and TensorFlow. By automating this mapping, Ember unlocks the full performance potential of the TMU without imposing an additional programmability burden.
\end{abstract}

\begin{abstract}
\textbf{CA}---Moltes càrregues de treball científiques i d’aprenentatge automàtic modernes presenten una estructura inherentment dispersa i irregular. Tanmateix, els superordinadors i centres de dades tradicionals estan optimitzats per a computació densa i regular. Com a resultat, les càrregues de treball disperses només utilitzen una fracció dels recursos disponibles del sistema, deixant sense aprofitar un rendiment significatiu. Tot i que s’han proposat solucions prèvies, aquestes solen adreçar-se a classes estretes d’operacions d’àlgebra de tensors dispersos, no aprofiten plenament els nuclis vectorials moderns, o deixen gran part de la càrrega de programabilitat a l’usuari, cosa que en limita l’adopció generalitzada. Per abordar aquesta mancança, aquesta tesi presenta un codisseny hardware--software que permet l’execució eficient d’operacions generals d’àlgebra de tensors dispersos a escala. Comencem analitzant les implicacions arquitectòniques de la dispersió en una àmplia gamma d’aplicacions científiques i d’aprenentatge automàtic. La nostra anàlisi revela que el recorregut de tensors dispersos i la fusió de tensors dispersos constitueixen els dos principals colls d’ampolla en arquitectures convencionals, sovint aturant l’execució fins al 99\% del temps. A partir d’aquestes observacions, introduïm la \textit{Tensor Marshaling Unit} (TMU), una unitat arquitectònica especialitzada que accelera aquestes operacions. Demostrem que la integració de la TMU en un processador CPU multinúcli millora el rendiment en un ordre de magnitud, superant les GPU d’última generació tant en rendiment com en eficiència energètica. Per obtenir aquests guanys de rendiment sense exposar la complexitat de la programació DAE als usuaris finals, també presentem Ember, un compilador que tradueix operacions disperses d’aprenentatge automàtic a codi \textit{Decoupled Access--Execute} (DAE) per a arquitectures com ara els sistemes TMU--CPU. Ember està implementat en MLIR, fet que permet una integració natural amb PyTorch i TensorFlow. Mitjançant l’automatització d’aquesta transformació, Ember desbloqueja tot el potencial de rendiment de la TMU sense imposar cap càrrega addicional de programabilitat.
\end{abstract}

\begin{abstract}
\textbf{ES}---Muchas cargas de trabajo científicas y de aprendizaje automático modernas presentan una estructura inherentemente dispersa e irregular. Sin embargo, los superordenadores y centros de datos tradicionales están optimizados para computación densa y regular. Como resultado, las cargas de trabajo dispersas utilizan solo una fracción de los recursos disponibles del sistema, dejando sin explotar un rendimiento significativo. Aunque se han propuesto soluciones previas, estas suelen dirigirse a clases estrechas de operaciones de álgebra de tensores dispersos, no aprovechan plenamente los núcleos vectoriales modernos, o dejan gran parte de la carga de programabilidad al usuario, lo que limita su adopción generalizada. Para abordar esta carencia, esta tesis presenta un codiseño hardware--software que permite la ejecución eficiente de operaciones generales de álgebra de tensores dispersos a escala. Comenzamos analizando las implicaciones arquitectónicas de la dispersión en una amplia gama de aplicaciones científicas y de aprendizaje automático. Nuestro análisis revela que el recorrido de tensores dispersos y la fusión de tensores dispersos constituyen los dos principales cuellos de botella en arquitecturas convencionales, llegando a bloquear la ejecución hasta el 99\% del tiempo. A partir de estos resultados, introducimos la \textit{Tensor Marshaling Unit} (TMU), una unidad arquitectónica especializada que acelera estas operaciones. Demostramos que la integración de la TMU en un procesador CPU multinúcleo mejora el rendimiento en un orden de magnitud, superando a las GPU de última generación tanto en rendimiento como en eficiencia energética. Para ofrecer estas mejoras de rendimiento sin exponer la complejidad de la programación DAE a los usuarios finales, también presentamos Ember, un compilador que traduce operaciones dispersas de aprendizaje automático a código \textit{Decoupled Access--Execute} (DAE) para arquitecturas como los sistemas TMU--CPU. Ember está implementado en MLIR, lo que permite una integración natural con PyTorch y TensorFlow. Al automatizar esta transformación, Ember desbloquea todo el potencial de rendimiento de la TMU sin imponer una carga adicional de programabilidad.
\end{abstract}
